\bookchapter{Bab 12: Hiburan Audio-Visual \& Pembacaan}

\chapterintro{Telefon sekarang asasnya telah menggantikan MP3, alat kawalan jauh TV, Kindle dan kons permainan tangan. Tapi sejujurnya, ramai orang hanya guna fungsi paling asas App ini, masih banyak kegunaan bagus yang belum ditemui. Bab ini mari kita cakap cara buat hiburan telefon lebih lancar, lebih selesa.}

\booksection{Dengar muzik: Pilih App mana, cara dengar macam mana baru selesa}

: Hak cipta paling lengkap, Jay Chou, Mayday dan artis popular pada asasnya hanya ada di QQ Music. Pilihan kualiti bunyi juga sangat kaya, ahli boleh buka kualiti bunyi tanpa kehilangan. Jika anda ambir berat tentang “lagu yang mahu dengar semua boleh cari”, QQ Music adalah pilihan paling stabil.

\begin{itemize}
  \item \textbf{Apple Music}: Pilihan utama pengguna Apple. Antara muka bersih, tiada iklan, kualiti bunyi baik, dan menyokong tanpa kehilangan dan Dolby Atmos. Tapi algoritma cadangannya tak sebaik App tempatan memahami citarasa pengguna Cina, budaya senarai lagu juga tak setara berkembang seperti NetEase Cloud. Sesuai untuk mereka yang suka pengalaman ringkas, tak ambil pusing fungsi sosial.
  \item \textbf{Spotify}: Pengguna global paling ramai, algoritma cadangan diakui paling kuat. "Discover Weekly" dan "Daily Mix" ketepatan tinggi sampai menakutkan, makin lama guna makin faham anda. Kekurangannya versi percuma banyak iklan, dan sesetengah lagu Cina hak cipta tak lengkap. Perlukan rangkaian khas baru boleh guna.
  \item \textbf{Beberapa tip dengar lagu kecil}: Jika anda guna fon kepala bluetooth, pergi ke tetapan App muzik tengok ada "Larasan Kesan Bunyi" atau "Equalizer", laraskan sikit bass dan treble, kesan peningkatan jelas. Tambahan pula, dengar lagu streaming sangat menggunakan data, ingat buka "Main tanpa kehilangan hanya dalam WiFi" dalam tetapan, atau muat turun senarai lagu kegemaran dahulu untuk dengar luar talian.
  \item \textbf{2. Podcast: "Rakan telinga" dalam perjalanan kerja}Podcast ini, orang yang tak dengar rasa tak berkaitan dengan diri, sekali mula dengar tak boleh kembali. Asasnya adalah "versi Netflix rancangan radio"—ada orang cakap teknologi, ada orang cerita sejarah, ada orang baca cerita, apa jenis pun ada.
\end{itemize}

\textbf{App Xiaoyuzhou}: App podcast terbaik dalam negara, antara muka bersih, tiada fungsi merumitkan. “Cadangan Harian” dan “Pilihan Editor” bantu anda temui rancangan baharu, ruang komen juga sangat aktif. Saya peribadi paling banyak dengar “Hu Zuo Hu You” (jenis budaya sejarah) dan “Bu He Shi Yi” (jenis isu sosial), setiap episod macam mendengar perbualan beberapa orang pintar.

\booksection{Podcast: “Rakan telinga” dalam perjalanan kerja}

: Terbina dalam iPhone, tak perlu muat turun App tambahan. Integrasi dengan CarPlay paling baik, semasa memandu sangat senang.

\begin{itemize}
  \item \textbf{Tip pendengaran}: Kebanyakan App podcast menyokong main 1.2x atau 1.5x, selepas biasa boleh jimat banyak masa. Saya biasanya 1.2x dengar bahasa Cina, 1.5x dengar bahasa Inggeris. Tambahan pula mesti buka "Muat turun automatik episod baharu" dan "Padam automatik selepas main", supaya telefon anda tak bertimbun fail audio yang telah didengar.
  \item \textbf{3. Tonton video: Cara betul guna pemain dan streaming}Banyak orang tak tahu, format yang boleh dimainkan di telefon jauh lebih banyak daripada yang disediakan dalam App streaming. Pemain tempatan yang bagus boleh selesaikan banyak masalah.
  \item \textbf{Pemain VLC}: Sumber terbuka percuma, hampir semua format boleh main—MKV, AVI, FLAC, malah imej cakera ISO. Jika anda muat turun satu video boleh main di komputer tapi tak boleh di telefon, 80\% masalah format, pasang VLC selesai. Ia juga menyokong sarung sarian luaran, letak fail sarung sari .srt dan video dalam folder yang sama, VLC akan automatik muatkan.
\end{itemize}

\booksection{Tonton video: Cara betul guna pemain dan streaming}

: Jika anda mahu pengalaman lebih halus, App berbayar ini berbaloi dibeli. Antara muka cantik, kawalan gerakan sangat lancar—kiri leret ke atas bawah laraskan kecerahan, kanan leret ke atas bawah laraskan volume, kiri kanan leret laju maju undur. Menyokong DTS dan kesan bunyi Dolby, tonton filem rasa kehadiran lebih kuat.

\begin{itemize}
  \item \textbf{Bilibili (B Station)}: Sudah bukan sekadar laman web ACG, sekarang sumber pembelajaran di atasnya terlalu banyak. Peperiksaan pengajian tinggi, pengaturcaraan, fotografi, memasak, apa-apa tutorial ada. Budaya peluru adalah jiwa B Station, semasa tonton video buka peluru, rasa macam ada sekumpulan orang teman anda tengok bersama. Tapi tonton video pembelajaran cadangan tutup peluru, jika tidak mudah terganggu.
  \item \textbf{Cast ke TV}: Skrin telefon sebesar mana pun tak selesa tengok seperti TV. Dalam WiFi yang sama, buka butang cast "TV" di penjuru kanan atas App video, boleh memancarkan gambar ke TV. Kebanyakan TV pintar dan projektor menyokong, tengok drama filem pengalaman jauh lebih baik. Pengguna iPhone juga boleh guna AirPlay terus cermin keseluruhan skrin.
  \item \textbf{4. Pembacaan e-buku: Boleh baca buku dengan baik di telefon}Ada orang kata baca buku di telefon sakit mata, mudah terganggu, ini memang betul. Tapi kelebihan pembacaan telefon juga jelas—bila-bila masa di mana sahaja, maksimumkan penggunaan masa pecahan. Saya dalam metro, beratur, sebelum tidur, guna telefon baca banyak buku.
  \item \textbf{WeChat Reading}: Buat masa ini App pembacaan pengalaman terbaik dalam negara. Pustaka buku sangat lengkap, buku terbitan dan novel web ada. Ia guna mod "tempoh baca tukar syiling buku", makin banyak baca boleh tukar buku percuma, sangat mesra pelajar. Fungsi "Garis Bawah" boleh tengok orang lain tandakan perenggan mana, kadang-kadang nampak ayat bagus yang ditanda orang lain, memang ada resonansi. Tapi "Papan Kedudukan Pembacaan" agak kompetitif, mudah jadi "baca untuk tempoh masa", cadangan tutup.
\end{itemize}

\booksection{Pembacaan e-buku: Boleh baca buku dengan baik di telefon}

: Jika anda ada peranti Kindle, Kindle App telefon boleh segerak kemajuan pembacaan. Kualiti susun atur e-buku asli Kindle jauh lebih baik daripada PDF cetak rompak, tapi kedai buku Kindle China telah ditutup, perlu tukar ke rantau luar negara.

\begin{itemize}
  \item \textbf{Apple Books}: Terbina dalam iPhone, sesuai baca EPUB dan PDF. Antara muka minimalis, animasi tukar halaman sangat selesa. Jika anda muat turun fail e-buku dari internet, terus "Kongsi ke Books" boleh import.
  \item \textbf{Cadangan tetapan pembacaan}: Tak kira App mana, mesti tetapkan latar belakang ke warna kuning hangat (mod mata jaga), besarkan font, jarakkan jarak baris. Baca pada waktu malam buka "Mod Gelap"—latar hitam teks putih, mata jauh lebih selesa. Jangan latar putih teks hitam font kecil baca terlalu lama dalam cahaya malap, keletihan mata adalah berkumpul.
  \item \textbf{5. Permainan telefon: Main seronok, tak rosakkan telefon}Sikap saya terhadap permainan telefon ialah: main boleh, tapi jangan biar ia sebaliknya kawal anda. Guna alat permainan secara wajar, main lebih selesa, juga tak ganggu hal penting.
  \item \textbf{Mod permainan/Ruang permainan}: Telefon utama sekarang semua ada fungsi ini. Selepas buka telefon akan automatik bersihkan latar belakang, sekatan pemberitahuan, halang sentuhan silap—anda sedang pertempuran kumpulan, tiba-tiba satu panggilan keluar, perkara seperti ini mod permainan boleh halang untuk anda. Sebahagian telefon juga menyokong mod prestasi/Mod E-sukan, paksa pemproses penuh, tapi bayarannya adalah haba dan bateri lebih cepat habis.
\end{itemize}

\booksection{Permainan telefon: Main seronok, tak rosakkan telefon}

: Jangan buta tarik maksimum. Pastikan kadar bingkai (kelembutan) dahulu, baru kejar kualiti gambar. 60 bingkai jauh lebih selesa daripada 30 bingkai, tapi 120 bingkai tekanan telefon sangat besar, telefon kelas pertengahan tak perlu paksa. Jika anda main permainan kompetitif (contohnya Honor of Kings, PUBG), keutamaan kadar bingkai jauh lebih tinggi daripada kualiti gambar—operasi lancar seratus kali lebih penting daripada rumput yang cantik.

\begin{itemize}
  \item \textbf{Telefon panas semasa main permainan macam mana}: Ini normal, tapi suhu tinggi jangka panjang akan rosakkan jangka hayat bateri. Beberapa cara sejuk: tanggalkan kes telefon untuk dissipasi; turunkan kecerahan skrin ke kira-kira 70\% (kecerahan adalah sumber haba tinggi); jangan cas sambil main permainan (haba bertindih, bateri cepat mati). Jika telefon panas sampai tak boleh pegang dengan tangan, cepat keluar biar ia rehat sekejap.
  \item \textbf{Alat kecil anti-ketagihan}: "Masa Skrin" iOS dan "Digital Wellbeing" Android boleh tetapkan had masa harian untuk App permainan. Saya tetapkan untuk diri hari kerja maksimum satu jam, hujung minggu dua jam. Sampai masa ikon App akan jadi kelabu, tak boleh klik masuk. Had ini betul-betul berguna, bukan sekadar hiasan.
  \item \textbf{Permainan awan}: Jika telefon anda prestasi tak cukup tapi mahu main permainan besar, boleh cuba platform permainan awan (contohnya Tencent START Cloud Gaming, NetEase Cloud Gaming). Permainan berjalan di pelayan jauh, gambar dihantar ke telefon anda dalam bentuk strim video. Keperluan rangkaian tinggi tapi hampir tiada keperluan terhadap prestasi telefon anda. Kekurangannya kelewatan tak boleh dielakkan, sesuai main strategi, RPG, tak sesuai FPS menembak.
  \item \textbf{Cuba buat: Bina sistem hiburan telefon anda}Dalam App muzik muat turun 3 album atau senarai lagu kegemaran anda, pastikan boleh dengar luar talian.
  \item \textbf{Pergi ke Xiaoyuzhou atau Apple Podcasts langgan 3 saluran yang menarik, muat turun satu episod untuk cuba dengar.}Pasang pemain VLC, cuba guna ia main satu fail video yang disalin dari komputer.
\end{itemize}

\begin{tcbpractice}{Dalam WeChat Reading cari satu buku yang mahu baca, tetapkan latar belakang warna hangat dan mod mata jaga, laraskan font ke saiz yang sesuai untuk anda.}
\begin{enumerate}
  \item Jika anda main permainan, buka mod permainan tengok ada tetapan apa yang boleh dioptimumkan (gangguan, keutamaan kadar bingkai dll.).
  \item Bab 13: Cadangan Aplikasi Alat
  \item Bab ini memilih aplikasi alat yang benar-benar boleh tingkatkan kecekapan, promosi kesihatan, dan memudahkan perjalanan.
  \item 1. Alat kecekapan: Urus masa \& perhatian anda
  \item Senarai tugasan
\end{enumerate}

\end{tcbpractice}
